% !TEX root = ../oddp.tex

\section{Alexander duality on the simplex and tensor products}\label{s:3complexes}

In this section we present a duality on the normalized chains of the augmented simplex. The main result of this section is \cref{prop:milnor_to_minimal}, where we find that building a connected $r$-cyclic diagonal is equivalent to build a map between two augmented resolutions of the cyclic group that intertwines certain endomorphisms of the resolutions. 

\begin{warning}
	Recall from Section \cref{s:simplices} and its three warnings that we are not following the usual convention to denote the simplex category.
\end{warning}
\subsection{Alexander duality on the augmented standard simplex}\label{s:alexander}

%Given a simplicial set $X$, let $\chains(X) = \sus{}\cadenas(X)$ be the right suspension of the normalized chains of $X$. \federico{I propose that }

\begin{definition}\label{d:poincare_duality_algebra}
	A connected commutative algebra $A$ of finite type is a \emph{Poincar\'e duality algebra} of \emph{formal dimension} $d$ if:
	\begin{enumerate}
		\item\label{i:pd1} $A_i = 0$ for $i > d$,
		\item\label{i:pd2} $\dim A_d = 1$,
		\item\label{i:pd3} $A_i \ot A_{d-i} \to A_d$ is non-degenerate.
	\end{enumerate}
\end{definition}

\begin{definition}\label{def:join_product}
	For any $[n] \in \ob\asimplex$ the \textit{join product} $\ast \colon \chains(\asimplex^{n})^{\ot 2} \to \chains(\asimplex^{n})$ is the linear map defined by sending a basis element $[v_0, \dots, v_{p-1}] \ot [v_{p},\dots,v_{m-1}]$ to $(-1)^{\sign{\perm}}[v_{\perm(0)}, \dots, v_{\perm(m-1)}]$	if $v_i \neq v_j$ for all $i \neq j$, where $\perm$ is the permutation ordering the vertices, and to $0$ otherwise. \federico{Here we can avoid the permutation $\pi$ writing $[v_0,\ldots,v_{p-1},v_p,\ldots,v_{m-1}]$}
\end{definition}

\begin{notation}
	If $\tau = [v_0,\dots,v_{k-1}]$ and $\tau'=[u_0,\dots,u_{m}-1]$ are two non-degenerate simplices of $\asimplex^n$, let $\lambda(\tau,\tau')$ be the parity of the permutation that orders the sequence $[v_0,\dots,v_{k-1},u_0,\dots,u_{m-1}]$. Let $\tau^c$ be the complementary face of $\tau$. The top generator $[0,1,\dots,n-1]$ will be denoted $\sigma$.
\end{notation}

\begin{theorem}
	The join product defines on $\chains(\asimplex^{n})$ the structure of a Poincar\'e duality algebra with unit the empty simplex $[0] \to [n]$ and formal dimension $n$. This structure is natural with respect to simplicial maps.
\end{theorem}

\begin{proof}
	The complex $\chains(\asimplex^{n})$ is connected and satisfies \cref{i:pd2,i:pd3} in \cref{d:poincare_duality_algebra} since $\chains(\asimplex^{n})_0 \cong \Z\{[0] \to [n]\}$, $\chains(\asimplex^{n})_{n} \cong \Z\{[n] \to [n]\}$, and $\chains(\asimplex^{n})_{n+k} \cong 0$ for $k>0$.

	That the join product is a natural chain map can be easily verified and a complete proof is presented in \cite[p.19]{medina2020prop1}.

	Thinking about the join product in terms of the union of sets with a permutation sign leads to a direct verification of its commutativity (in the graded sense) and unitality with respect to the empty simplex.

	To verify \cref{i:pd3} consider a basis element $x = [v_1,\dots,v_i]$.
	Let $x^c$ be the ordered complement of $\set{v_1,\dots,v_i}$ in $\{0,\dots,n-1\}$ and notice that $x \ast x^c = \pm [0,\dots,n-1]$ as required.
\end{proof}

This pairing thus induces an isomorphism
\begin{equation}\label{eq:iso1}
	\alex \colon \chains(\asimplex^n) \to \susp{n}\cochains(\asimplex^n)
\end{equation}
that sends a chain $\tau$ to $(-1)^{\lambda(\tau,\tau^c)}(\tau^c)^\vee$. We will refer to this isomorphism as \emph{Alexander duality}.

\begin{remark}\label{remark:alex}
	If $\tau = [v_0,\dots,v_{k-1}]$ and $\tau^c = [u_0,\dots,u_{m-1}]$ are the ordered representatives, then
	\begin{align*}
		\lambda(\tau,\tau^c)&\equiv \sum_{i=0}^{k-1} (v_i-i) \equiv \sum_{i=0}^{k-1} v_i - \binom{k+1}{2} \\
		\lambda(\tau,\tau^c)&\equiv \sum_{j=0}^{m-1} (n-u_j-m+j) \equiv \sum_{j=0}^{m-1} u_j +(n-m)(m+1)+\binom{m+1}{2}.
	\end{align*}
\end{remark}

\begin{notation}
	We will denote normalized chains between brackets and normalized cochains between parentheses.
\end{notation}

\subsection{Simplicial functor tensor products}\label{s:functortensor}

Recall that if $A$ is an augmented cosimplicial chain complex and $B$ is an augmented simplicial chain complex, the functor tensor product $A \ot_{\asimplex} B$ is the quotient of the tensor product $A \ot B$ by the relation generated by
\begin{align*}
	\d^i a \ot b &\sim (-1)^{|\d^i||a|}a \ot \d_i b
	&
	\s^ia \ot b &\sim (-1)^{|\s^i||a|}a \ot \s_ib.
\end{align*}

The chain complexes $\chains(\asimplex^\bullet)$ form an augmented cosimplicial chain complex with
\begin{align*}
	\d^i \colon \chains(\asimplex^n)& \lra \chains(\asimplex^{n+1})
	&
	\s^i \colon \chains(\asimplex^{n})& \lra \chains(\asimplex^{n-1})
\end{align*}
given on an ordered representative $[v_0,\dots,v_{k-1}]\in \chains(\asimplex^n)$ by
\begin{align*}
\d^i([v_0,\dots,v_{k-1}]) &= [v_0,\dots,v_{j-1},v_j+1,\dots,v_{k-1}+1] \quad \text{if $v_{j-1}<i\leq v_j$}
\\
\s^i([v_0,\dots,v_{k-1}]) &= [v_0,\dots,v_{j-1},v_j-1,\dots,v_{k-1}-1] \quad \text{if $v_{j-1} \leq i < v_j$}.
\end{align*}
If $X_\bullet$ is an augmented simplicial $R$-module, the chain complex $\chains(X)$ can alternatively be described as the functor tensor product over $\asimplex$ of the augmented cosimplicial chain complex $\chains(\asimplex^{\bullet})$ and $X_\bullet$:
\begin{equation}\label{eq:1} 
 \chains(\asimplex^\bullet) \ot_{\asimplex} X_\bullet
 \cong	
 \chains(X).
\end{equation}
%This sends a simplex $[k] \to [n]$ and a simplex $[n] \to X$ to the composition $[k] \to X$.



\subsection{Alexander duality}
Via Alexander duality, the sequence of chain complexes $\susp{\bullet }\cochains(\asimplex^\bullet)$ is then endowed with the following augmented \emph{cosimplicial} structure
\begin{align*}
	\d^i \colon \cochains(\asimplex^n) & \lra \susp{}\cochains(\asimplex^{n+1})
	\\
	\s^i \colon \cochains(\asimplex^{n+1}) & \lra \susp{-1}\cochains(\asimplex^n)
\end{align*}
given on an ordered generator $U=(u_0,\dots,u_{m-1})$ by
\begin{align*}
	\d^i(U) &=
	(-1)^{i+j+m+n+1}(u_0,\dots,u_{j-1},i,u_{j}+1,\dots,u_{m-1}+1)
	\quad \text{if $u_{j-1}<i \leq u_{j}$}
	\\
	\s^i(U) &= \begin{cases}
	(-1)^{i+j+m+n+1}(u_0,\dots,u_{j-1},u_{j+1}-1\dots,u_{m-1}-1) & \text{if $u_{j} = i$,}
	\\
	0 & \text{otherwise.}
	\end{cases}
\end{align*} %Why is the next statement important? It wasn't stated for $\chains(\asimplex^\bullet)$.
 The chain complexes $\cochains(\asimplex^\bullet)$ have a simplicial structure as well: the linear dual of the cosimplicial structure on $\chains(\asimplex^\bullet)$, whose face and degeneracies are denoted by $\d_i$ and $\s_i$.
 
The following isomorphism gives another characterization of $\chains(X)$:
\begin{equation}\label{eq:dual}
	\alex^{-1} \ot \Id\colon \susp{\bullet }\cochains(\asimplex^\bullet)\ot_{\asimplex} X_\bullet
	\lra
	\chains(\asimplex^{\bullet})\ot_{\asimplex} X_\bullet
\end{equation}
The $r$-fold tensor product of (co)simplicial chain complexes is again a (co)simplicial chain complex, and the functor tensor product construction commutes with the tensor product:
\begin{equation}\label{eq:assemble}
	\susp{\bullet r}\cochains(\asimplex^\bullet)^{\ot r}\ot_{\asimplex} X^{\ot r}_\bullet 
	\lra 
	\left(\susp{\bullet }\cochains(\asimplex^\bullet)\ot_{\asimplex} X_\bullet \right)^{\ot r}.
\end{equation}

\subsection{Simplicial tensor products and functor tensor products}
A generator $U_0\ot\cdots \ot U_{r-1}$ of $\cochains(\asimplex^n)^{\ot r}$ is \emph{almost $i$-full} if all but one of the tuples $U_0,\ldots,U_{r-1}$ contain the number $i$. It is \emph{$i$-full} if all the tuples contain $i$. The generators that are not $i$-full for any $i$ are also known as \emph{interior} for the cosimplicial structure. Let $\cochains(\asimplex^n)^{\ot r}_{\interior}$ denote the quotient of the complex $\cochains(\asimplex^n)^{\ot r}$ by the subcomplex generated by the non-interior generators, and let $\chains(\asimplex^n)^{\ot r}_\interior$ be its linear dual. The cyclic group $\Cyc_r$ acts on them permuting the factors.
\begin{definition}
	The sequence of $R\Cyc_r$-chain complexes $[n]\mapsto \susp{n(r-1)}\cochains(\asimplex^{n})^{\ot r}_\interior$ has a semi-simplicial structure with face maps 
	\begin{align*}
		\d_i\colon \cochains(\asimplex^n)^{\ot r}_\interior &\lra
		\susp{1-r}\cochains(\asimplex^{n-1})^{\ot r}_\interior	
	\end{align*}
	given by $\d_i = \sum \s^i\ot \ldots \s^i\ot\d_i\ot \s^i \ot \ldots \s^i$. This map sends an almost $i$-full generator $U_0\ot \cdots \ot U_{r-1}$ to $\pm \s^iU_0\ot\ldots\ot \d_iU_j\ot\ldots\ot \s^iU_{r-1}$ if $i$ is not contained in $U_j$. All other generators are sent to $0$. Dualising we obtain the cosemi-simplicial chain complex $\sus{\bullet(1-r)}\chains(\asimplex^\bullet)^{\ot r}_\interior$ whose coface maps are denoted $\d^i$.	
\end{definition}

 For the following lemma recall that $\ot$ denotes the semi-simplicial tensor product of \cref{s:simplices} and $\ot_{\asimplex}$ denotes the functor tensor product of \cref{s:functortensor}.

\begin{lemma}\label{lemma:simp-cosimp}
	The map 
	\[\chains(\susp{\bullet(r-1)}\cochains(\asimplex^\bullet)^{\ot r}_\interior\ot (X^\nd)^{\ot r}_\bullet)
	\lra \susp{\bullet r}\cochains(\asimplex^\bullet)^{\ot r} \ot_{\asimplex} X ^{\ot r}_\bullet\]
	that sends a generator $(U_0\ot\ldots \ot U_{r-1})\ot (x_0\ot\ldots\ot x_{r-1})$ to the same generator on the right hand-side is an isomorphism of $R\Cyc_r$-chain complexes.
\end{lemma}
\begin{proof}
	TBD. The paper \url{https://www.jstor.org/stable/2043689} basically says that this is an isomorphism of graded modules (i.e., forgetting the differential). Probably this is already in the literature.
\end{proof}

\subsection{A connected diagonal}

Recall from \cref{s:preliminaries} the chain complex $\Wdualaug{r}$ and the chain map $\theta^\vee\colon \Wdualaug{r}\to \susp{1-r}\Wdualaug{r}$.

\begin{definition} The sequence of chain complexes $[n]\mapsto \susp{n(r-1)}\Wdualaug{r}$ has a semi-simplicial structure whose face maps are all equal to $\theta^\vee$. Dualising we obtain the cosemi-simplicial chain complex $\sus{\bullet(1-r)}\Waug{r}$ whose coface maps are all equal to $\theta$ (cf. \cref{ex:311} and \cref{lemma:adj}).
\end{definition}

\begin{lemma}\label{lemma:hotimessimp}
	The chain complex $\susp{\bs(r-1)}\Wdualaug{r}\ot \chains(X)_\bs$ is isomorphic to the chain complex $\chains(\susp{\bullet(r-1)}\Wdualaug{r} \ot X^\nd_\bullet)$
\end{lemma}
\begin{proof}
	This is the isomorphism \eqref{eq:adj}.
\end{proof}
Let $\chains(\asimplex^\infty)^{\ot r}_{\mathrm{int}}$ be the colimit of the maps $\chains(\asimplex^n)^{\ot r}_{\mathrm{int}} \to \chains(\asimplex^{n+1})^{\ot r}_{\mathrm{int}}$ induced by the coface maps
$\d^i\colon \chains(\asimplex^n)^{\ot r} \to \chains(\asimplex^{n+1})^{\ot r}$.


\begin{corollary}
	A family of equivariant chain maps 
	\[
	\Psi_n\colon \chains(\asimplex^{n})^{\ot r}_\interior \lra \Waug{r}
	\]
	such that $\theta\circ \Psi_n =  \Psi_{n+1} \circ \d^i$ for each $i$ yields an equivariant chain map
	\[
	\mu\colon \susp{\bs(r-1)}\Wdualaug{r}\ot \chains(X)_\bs\lra \chains(X)^{\ot r}.
	\]
	If additionally $\Psi_n$ factors as $c^n\cdot \Psi^{\infty}\circ \iota_n$
	\[
		\chains(\asimplex^n)^{\ot r}_{\mathrm{\mathrm{int}}}\overset{\iota_n}{\lra} \chains(\asimplex^\infty)^{\ot r}_{\mathrm{int}} \overset{\Psi^{\infty}}{\lra} \Waug{r}
	\]
	for some cosemi-simplicial map $\Psi^{\infty}$ and some constant $c$,
	then they satisfy that 
	\begin{align*}
		\mu\circ \sus{} &= c\cdot \sus{\ot r}\circ \mu
		\\
		\mu\circ \susp{} &= c\cdot \rho(\susp{\ot r}\circ \mu).
	\end{align*}
\end{corollary}
\begin{proof}
	The map is obtained as the composition of the isomorphism of Lemma \ref{lemma:hotimessimp}, the map \martin{change $n$ by $\bullet$?}
	\[
		\Psi_n^\vee\ot \diag\colon \chains(\susp{\bullet(r-1)}\Wdualaug{r} \ot X^{\nd}_\bullet) 
		\lra
		\chains\left(\susp{\bullet(r-1)}\cochains(\asimplex^n)^{\ot r}_\interior\ot (X^{\nd})^{\ot r}_\bullet\right),	
	\]
	the isomorphism of Lemma \ref{lemma:simp-cosimp}, the isomorphism \eqref{eq:assemble}, the $r$-fold tensor product of the isomorphism \eqref{eq:dual} and the $r$-fold tensor product of the isomorphism \eqref{eq:1}.
\begin{tikzpicture}
	\node (1) at (-3.8,3){$\susp{\bs(r-1)}\Wdualaug{r}\ot \chains(X)_\bs$};
	\node (2) at (-3.8,0){$\chains\left (\susp{\bullet(r-1)}\Wdualaug{r}\ot X^\nd_\bullet\right )$};
	\node (3) at (-3.8,-3){$\chains\left (\susp{\bullet(r-1)}\cochains(\asimplex^\bullet)^{\ot r}_\interior\ot (X^\nd)^{\ot r}_\bullet \right )$};
	\node (4) at (3.5,-3){$\susp{\bullet r}\cochains(\asimplex^\bullet)^{\ot r}\ot_{\asimplex} X^{\ot r}_\bullet $};
	\node (5) at (3.5,-1){$\left (\susp{\bullet}\cochains(\asimplex^\bullet)\ot_{\asimplex}X_\bullet\right)^{\ot r}$};
	\node (6) at (3.5,1){$\left (\chains(\asimplex^\bullet)\ot_{\asimplex}X_\bullet\right )^{\ot r}$};
	\node (7) at (3.5,3){$\chains(X)^{\ot r}$};
	%
	\draw[->] (1)--node[above]{$\mu$}(7);
	\draw[->](1)--node[right]{\ref{lemma:hotimessimp}}node[left]{$\cong$}(2);
	\draw[->](2)--node[right]{$\Psi^\vee\ot \diag$}(3);
	\draw[->](3)--node[above]{\ref{lemma:simp-cosimp}}node[below]{$\cong$}(4);
	\draw[->](4)--node[right]{\eqref{eq:assemble}}node[left]{$\cong$}(5);
	\draw[->](5)--node[right]{$(\Lambda^{-1} \ot \Id)^{\ot r}$}node[left]{$\cong$}(6);
	\draw[->](6)--node[right]{$\eqref{eq:1}^{\ot r}$}node[left]{$\cong$}(7);
\end{tikzpicture}
	
	THE SECOND PART (SUSPENSION) IS TBD AND POSSIBLY SLIGHTLY ADJUSTED.
\end{proof}

In Sections \ref{s:milnor}, \ref{s:milnor_to_per} and \ref{s:per_to_minimal} we will define two additional cosemi-simplicial chain complexes $	\sus{\bullet(1-r) }\chains(\partial \asimplex^{r})^{\ot \bullet}$ and $\sus{\bullet(1-r)}\Per{r}$ with an action of $\Cyc_r$ and in \cref{lemma:bartomilnor}, \cref{prop:milnor_to_per} and \cref{prop:per_to_minimal} we will construct three cosemi-simplicial maps between them (for the sake of clearness, we omit the symbol $\sus{\bullet(1-r)}$ everywhere):  
\[
	\chains(\asimplex^{\bullet})^{\ot r}_{\interior}
	\overset{\beta\circ \alpha}{\lra}
	\chains(\partial \asimplex^{r})^{\ot \bullet }
	\overset{\psi}{\lra}
	\Per{r}
	\overset{\phi}{\lra}
	\Waug{r}
\]
whose composition satisfies the conditions of this corollary. There is a map $\psi$ for each choice of $r$-cyclic straightening with duality. The composition $\phi\circ\psi\circ\beta\circ\alpha$ satisfies the extra hypothesis of the Corollary.
